% BEGIN REV09_FRICKE_JET
\subsection{Correspondance de Fricke et récupération différentielle}
\label{corpus9:iv:fricke}

La trace réticulaire de niveau deux et sa complétion de Fricke
fournissent, après la construction exceptionnelle, les
coordonnées
\[
 t=t_2,\qquad s=\frac{4096}{t},\qquad
 T=t+s+24,\qquad t\in\mathbb C^\times.
\]
L’involution échange \(t\) et \(s\). Les deux lectures modulaires
ordonnées sont
\[
 j_1=\frac{(t+256)^3}{t^2},\qquad
 j_2=\frac{(s+256)^3}{s^2}.
\]
Ces formules utilisent l’uniformisation déjà établie dans
\eqref{km:j-nivel2} et conservent le choix de la lecture orientée.

\begin{proposition}[Correspondance et discriminant]
\label{corpus9:iv:fricke-correspondencia}
On a les identités
\begin{align}
 j_1-j_2&=(s-t)(T+23),\\
 j_1+j_2&=T^2+T-7256,\\
 j_1j_2&=(T+248)^3.
 \label{corpus9:iv:fricke-simetricos}
\end{align}
Les deux valeurs sont racines de
\[
 X^2-(T^2+T-7256)X+(T+248)^3=0,
\]
dont le discriminant est
\begin{equation}
 \Delta=(T-152)(T+104)(T+23)^2.
 \label{corpus9:iv:fricke-discriminante}
\end{equation}
\end{proposition}
\begin{proof}
Le développement de la première lecture, en utilisant \(ts=4096\), donne
\(j_1=t+768+48s+s^2\) ; la seconde s’obtient en échangeant
\(s,t\). La différence donne \((s-t)(t+s+47)\), et la somme se réduit à
\((t+s)^2+49(t+s)-6656\). La substitution de \(t+s=T-24\) produit
les deux premières identités. Pour le produit,
\[
 j_1j_2
 =\frac{((t+256)(s+256))^3}{(ts)^2}
 =(T+248)^3.
\]
Enfin,
\((s-t)^2=(T-24)^2-16384=(T-152)(T+104)\).
Multiplier par \((T+23)^2\) prouve
\eqref{corpus9:iv:fricke-discriminante}.
\end{proof}

\begin{proposition}[Récupération de l’orientation par les valeurs et le premier jet]
\label{corpus9:iv:fricke-jet}
Pour \(T\ne-23\), la paire ordonnée \((T,j_1)\) récupère
\begin{equation}
 t=\frac{T^2-j_1-3904}{T+23}.
 \label{corpus9:iv:fricke-inversa}
\end{equation}
Dans \(T=-23\), les deux paramètres

\[
 t_\pm=\frac{-47\pm45i\sqrt7}{2}
\]
partagent \((T,j_1)=(-23,-3375)\). Le premier jet
\(p=dj_1/dT\), calculé sur la branche paramétrée par \(t\),
les sépare et permet de récupérer
\begin{equation}
 p_\pm=\frac{-45\mp45i\sqrt7}{2},\qquad
 s=p-1,\qquad t=-46-p.
 \label{corpus9:iv:fricke-jet-inversa}
\end{equation}
\end{proposition}
\begin{proof}
De \(s^2=(T-24)s-4096\) résulte

\[
 j_1=(T+23)s+T-3352
     =T^2-3904-(T+23)t.
\]
Pour \(T\ne-23\), cette égalité donne
\eqref{corpus9:iv:fricke-inversa}. Pour \(T=-23\),
les conditions \(t+s=-47\), \(ts=4096\) donnent
\(t^2+47t+4096=0\), de discriminant \(-45^2\cdot7\).
La même identité donne \(j_1=-3375\) aux deux racines.


La dérivée \(dT/dt=1-4096/t^2\) ne s’annule qu’en
\(t=\pm64\). Aucune des deux racines exceptionnelles n’a cette
valeur, de sorte que \(T\) est une coordonnée locale sur les deux branches.
En différentiant la première expression de \(j_1\),
\[
 \frac{dj_1}{dt}
  =(s+1)\frac{dT}{dt}+(T+23)\frac{ds}{dt}.
\]
Dans \(T=-23\), il vient \(p=s+1=-46-t\), qui prouve
\eqref{corpus9:iv:fricke-jet-inversa}, avec des signes opposés
pour les indices appariés \(t_\pm,p_\pm\).
\end{proof}

Les points fixes de Fricke, \(t=\pm64\), expriment respectivement
\((T,j)=(152,8000)\) et \((-104,1728)\). La coïncidence en
\(T=-23\) correspond en revanche à deux points distincts de la
paramétrisation. Le corps de définition s’obtient à partir de l’équation :
\[
 \mathbb Q(t_\pm)=\mathbb Q(\sqrt{-7}),\qquad
 \left(\frac{2t+47}{45}\right)^2=-7.
\]
Sur cette fibre, Fricke coïncide avec la conjugaison quadratique,
de trace \(-47\) et de norme \(4096\). Le jet conserve ce corps,
car \((2p+45)^2=-45^2\cdot7\). La reconstruction est exacte
dans cette correspondance modulaire : la donnée différentielle distingue
les deux orientations qui partagent leurs valeurs scalaires.
% END REV09_FRICKE_JET
